\chapter{Introduzione}
Negli ultimi anni l’evoluzione dei Large Language Models (LLM) ha reso possibile l’utilizzo di sistemi di intelligenza artificiale in grado di assistere attivamente il processo di sviluppo software. Strumenti come GitHub Copilot, ChatGPT, CodeWhisperer e altri sistemi basati su modelli linguistici sono oggi in grado di generare codice, correggere errori e supportare la progettazione di componenti software a partire da una descrizione del problema.
\\
Queste capacità stanno progressivamente modificando il modo in cui gli sviluppatori affrontano il processo di programmazione. Tradizionalmente, lo sviluppo software consiste nella progettazione di un algoritmo e nella sua traduzione in un linguaggio di programmazione formale. Con l’introduzione dei modelli linguistici, invece, parte di questo processo può essere delegata a sistemi automatici capaci di generare codice a partire da descrizioni ad alto livello. L’utilizzo dell’intelligenza artificiale nella programmazione software introduce la necessità di ripensare il workflow di sviluppo a partire dalla analisi del problema fino al debugging e modifica del codice.
\\
In questo contesto, uno degli aspetti centrali diventa la capacità di comunicare efficacemente con il modello linguistico, fornendo descrizioni chiare del problema, del contesto e dei vincoli richiesti dalla soluzione. Tuttavia, l'interazione con questi sistemi non può essere ridotta alla semplice formulazione di una richiesta testuale. Per ottenere risultati affidabili è necessario integrare l'intelligenza artificiale all'interno di un workflow di sviluppo strutturato, in cui il modello possa accedere alle informazioni rilevanti del progetto e operare all'interno di un contesto ben definito.
\\ 
\section{Obiettivo}
Per analizzare concretamente queste dinamiche, il presente lavoro utilizza come caso di studio lo sviluppo reale di un'applicazione web commissionata per l'accettazione e la gestione degli ordini di una pizzeria. La realizzazione di questo sistema informatico non rappresenta quindi solamente l'obiettivo applicativo del progetto, ma costituisce anche il contesto sperimentale attraverso cui osservare e valutare il ruolo degli strumenti di intelligenza artificiale nello sviluppo del software.
\\
Nei seguenti capitoli verranno illustrate diverse modalità di utilizzo dei sistemi di intelligenza artificiale durante il processo di sviluppo, analizzando in particolare come questi strumenti possano supportare lo sviluppatore nelle varie fasi della progettazione e dell'implementazione del sistema. Verranno inoltre presentati diversi casi di ottimizzazione del processo di sviluppo ottenuti attraverso l'utilizzo di ambienti di lavoro controllati, in cui l'interazione tra sviluppatore e modello linguistico è stata strutturata per massimizzare l'efficacia del supporto fornito dall'intelligenza artificiale.
\\
Attraverso questo caso di studio sarà quindi possibile osservare come l'integrazione di strumenti basati su modelli linguistici possa influenzare il processo di sviluppo software, evidenziando sia i vantaggi operativi sia le criticità legate all'utilizzo di questi sistemi. L'obiettivo non è solamente dimostrare la capacità dei modelli di generare codice, ma analizzare in che modo tali strumenti possano essere utilizzati come supporto efficace all'interno di un processo di sviluppo strutturato.


\section{Strumenti Utilizzati}


Durante lo sviluppo dell'applicazione sono stati utilizzati diversi strumenti basati su modelli di intelligenza artificiale, integrati all'interno dell'ambiente di sviluppo con l'obiettivo di supportare le varie fasi del processo di programmazione. Questi strumenti sono stati impiegati per assistere lo sviluppatore nella progettazione dell'interfaccia utente, nella generazione e modifica del codice e nell'organizzazione della struttura del progetto.
\paragraph{GitHub Copilot}
GitHub Copilot è un assistente di programmazione basato su intelligenza artificiale integrato direttamente nell'ambiente di sviluppo Visual Studio Code. Il sistema è stato utilizzato principalmente per fornire suggerimenti automatici di codice durante la scrittura dei file sorgente.
Copilot analizza il contesto del codice e propone automaticamente completamenti e implementazioni di funzioni, contribuendo a velocizzare la scrittura del codice e ridurre operazioni ripetitive durante lo sviluppo.

\paragraph{Google Stitch}

Google Stitch è uno strumento sperimentale di progettazione di interfacce utente basato su intelligenza artificiale sviluppato da Google Labs. Il sistema permette di generare prototipi di interfacce web e mobile a partire da descrizioni testuali o da semplici wireframe.
Nel contesto del progetto, Stitch è stato utilizzato come strumento di supporto alla fase iniziale di progettazione dell'interfaccia grafica dell'applicazione. Attraverso l'utilizzo di descrizioni testuali è stato possibile generare rapidamente layout preliminari dell'interfaccia utente, successivamente adattati e integrati nel progetto finale.

\paragraph{Windsurf}
Windsurf è un ambiente di sviluppo integrato (IDE) progettato per integrare profondamente l'intelligenza artificiale nel processo di programmazione. A differenza dei tradizionali editor di codice, Windsurf è stato progettato come un editor ``AI-native'', in cui i modelli linguistici partecipano attivamente alla generazione, modifica e debugging del codice.
Questo strumento consente di interagire con l'intera codebase attraverso comandi in linguaggio naturale, permettendo al sistema di comprendere il contesto del progetto e di generare modifiche su più file in modo coordinato. Durante lo sviluppo del progetto, Windsurf è stato utilizzato per sperimentare modalità di sviluppo assistito da AI più avanzate rispetto ai tradizionali sistemi di completamento automatico del codice.

\section{Premessa: contesto e scalabilità del caso di studio}

Il sistema software analizzato in questo lavoro consiste in un'applicazione web per la gestione degli ordini e dell'organizzazione operativa di una pizzeria. L'applicazione è progettata come sistema informativo per attività di ristorazione di piccole-medie dimensioni e presenta una struttura architetturale centralizzata, basata su una singola istanza server e su un database locale SQLite. L'interfaccia utente è accessibile tramite browser web ed è pensata per essere utilizzata da un numero limitato di operatori simultanei.

È importante sottolineare che il dominio applicativo analizzato presenta un livello di complessità architetturale relativamente contenuto. Il sistema non include infatti alcune delle caratteristiche tipiche delle applicazioni enterprise di larga scala, come architetture distribuite basate su microservizi, sistemi di messaging asincrono, gestione avanzata della concorrenza o infrastrutture di scalabilità orizzontale. Analogamente, il numero limitato di utenti simultanei e la natura prevedibile dei dati trattati rendono il modello relativamente semplice.

La scelta di un sistema con complessità moderata consente di analizzare in modo più chiaro le modalità di integrazione degli strumenti basati su intelligenza artificiale nel processo di sviluppo software. In un contesto fortemente complesso, infatti, numerosi fattori architetturali e infrastrutturali potrebbero rendere più difficile isolare e osservare gli effetti dell'utilizzo dei modelli linguistici nel workflow di programmazione.

Da notare come i pattern di sviluppo assistito da intelligenza artificiale analizzati in questo lavoro non dipendono strettamente dal dominio applicativo specifico. Tali approcci risultano infatti trasferibili anche a sistemi software di maggiore complessità, nei quali le stesse tecniche possono essere applicate su architetture più articolate.

Il caso di studio presentato in questa tesi si colloca quindi in una posizione intermedia: sufficientemente complesso da presentare problemi reali di sviluppo software, ma al tempo stesso abbastanza contenuto da permettere un'analisi completa del processo di sviluppo assistito da intelligenza artificiale.

\section{Requisiti pratici del progetto}

L’applicazione sviluppata in questo lavoro nasce da una richiesta concreta da parte del proprietario di una pizzeria, il quale necessitava di uno strumento informatico in grado di supportare sia la gestione operativa del ristorante durante il servizio sia l’analisi economica dell’attività. L’obiettivo principale era disporre di un sistema semplice da utilizzare dal personale del locale, accessibile tramite browser e in grado di centralizzare la gestione degli ordini e delle informazioni economiche.

Oltre agli aspetti operativi legati al servizio, il proprietario era sopratutto interessato alla possibilità di analizzare i dati generati dall’attività del ristorante in maniera semi automatica, in particolare, il sistema doveva essere in grado di raccogliere informazioni sugli ordini effettuati e trasformarle in indicatori utili alla gestione economica del locale. Tra le funzionalità richieste rientravano la possibilità di monitorare i ricavi generati dal ristorante, analizzare il numero di clienti serviti in determinati intervalli temporali e valutare l’andamento dei prezzi e dei guadagni nel tempo.

Un ulteriore aspetto di interesse riguardava la possibilità di ottenere statistiche aggregate sull’attività del ristorante. Il proprietario desiderava, ad esempio, poter stimare il guadagno medio per coperto, osservare il numero medio di persone servite durante il servizio e individuare eventuali variazioni nei ricavi o nella domanda dei prodotti. L’obiettivo era quello di utilizzare i dati raccolti dal sistema come supporto alle decisioni gestionali, permettendo una visione più chiara dell’andamento economico del locale.